\section{Score Pretraining and Joint Receiver Training}
\label{sec:learning}

We first learn the channel score from perturbed channel samples and
then train the unfolded receiver on received pilots, jointly adapting
the channel denoisers, activity classifiers, and message parameters.

\subsection{First-Order Score Matching}

For a normalized channel sample $h_R\in\mathbb R^{2D}$, we form
the perturbed observation $u=h_R+\sqrt{v/2}\,\varepsilon$, where
$\varepsilon\sim\mathcal N(0,\I_{2D})$. Denote the perturbed
density, the real-coordinate form of $q_1$ in
\eqref{eq:twc-local-densities}, by $q_{1,R}(u;v)$ and its score by
$s^\star(u,v)=\nabla_u\log q_{1,R}(u;v)$.
Tweedie's formula relates this score to the conditional channel mean
\cite{efron2011,cai2026}:
\begin{equation}
\E[h_R\mid u]=u+\frac v2\,s^\star(u,v).
 \label{eq:twc-tweedie}
\end{equation}
Replacing $s^\star$ by a score network gives the denoiser in
\eqref{eq:twc-conditional-mean}, initially with $a_\ell=1$, without
direct integration over the channel distribution. To learn the
unknown marginal score, denoising score matching uses the Gaussian
perturbation score $-\varepsilon/\sqrt{v/2}$ as the training target
while preserving the score-matching minimizers
\cite{vincent2011,song2019}. With variance weighting, the loss is
\begin{equation}
 \mathcal L_{\rm DSM}(\theta)=\E_{h_R,v,\varepsilon}
 \left[\frac{\|\sqrt{v/2}\,s_\theta(h_R+\sqrt{v/2}\varepsilon,v)
                      +\varepsilon\|^2}{2D}\right].
 \label{eq:twc-dsm}
\end{equation}
Samples are drawn before activity masking and perturbed over a range
of noise levels, with continuous delays and Doppler shifts represented
in their effective-channel coefficients. The learned score supplies
the conditional channel mean under the active hypothesis; activity
probabilities are learned from received pilots in the next stage.

\subsection{Joint Channel and Activity Training}
\label{sec:twc-joint-training}

To adapt the pretrained denoisers to their layer inputs, each
training observation passes through the updates in
Section~\ref{sec:algorithm}, using the sensing matrix and noise variance
supplied by the physical model. The trainable parameters are
$\theta_\ell$, $\phi_\ell$, $c_{\alpha,\ell}$, $c_{\gamma,\ell}$,
and, for layers that return a message, $c_{\beta,\ell}$.
For a fixed distribution of denoiser inputs, expected squared error on
active users and expected binary cross-entropy on all users are minimized by the
conditional channel mean and activity probability, respectively
\cite{menon2016}, matching the two targets in
\eqref{eq:twc-mixture-mean}.

Suppose the receiver training set contains $B$ scenes, with $K_i$
candidate users in scene $i$. Let $h_{ik,R}$ and $\alpha_{ik}$
denote the channel and activity label of user $k$ in scene $i$.
Training and validation use disjoint physical scenes.
The channel and activity losses at layer $\ell$ are respectively
given by
\begin{align}
 R_H^{(\ell)}&=\frac{1}{E_{\rm tr}}\sum_{i,k}
       \alpha_{ik}\|\mu_{ik\ell}-h_{ik,R}\|^2,
 \label{eq:twc-channel-loss}\\
 R_a^{(\ell)}&=-\frac{1}{T_{\rm tr}}\sum_{i,k}
       \left[\alpha_{ik}\log p_{ik\ell}
       +(1-\alpha_{ik})\log(1-p_{ik\ell})\right].
 \label{eq:twc-activity-loss}
\end{align}
Here $E_{\rm tr}=\sum_{i,k}\alpha_{ik}\|h_{ik,R}\|^2$ and
$T_{\rm tr}=\sum_i K_i$ are the training set's active-channel energy
and user count, which are held fixed so that scenes with different
numbers of users retain the same normalization. The two normalized
losses receive equal coefficients in both training and model selection.

For a receiver with $m$ layers, the joint training objective is
\begin{equation}
 J_m=\frac{\sum_{\ell=1}^{m}\omega_\ell
                  \left(R_H^{(\ell)}+R_a^{(\ell)}\right)}
                 {\sum_{\ell=1}^{m}\omega_\ell},
 \label{eq:twc-layer-objective}
\end{equation}
where $\omega_\ell$ weights the loss at layer $\ell$, including
intermediate estimates used to form subsequent messages. Since the
classifier takes $\mu_{k\ell}$ as an input, gradients from the activity
loss update the channel denoiser within the same layer. Across layers,
the product $p_{k\ell}\mu_{k\ell}$ determines the residual and return
message, through which later channel losses update earlier activity
classifiers and message parameters. Joint training thus accounts for
both the current estimation errors and their effects on later layers.

\subsection{Progressive Training}

Each newly added layer has independent parameters and receives the
messages produced by the preceding layers. Its score network is
initialized from the same pretrained network, while its log-odds
correction $g_{\phi_\ell}$ and $c_{\alpha,\ell}$, $c_{\beta,\ell}$,
and $c_{\gamma,\ell}$ start at zero, giving an initial activity
probability of $\rho$.
Earlier layers retain their selected parameters.

At depth $m$, we first train the new layer with the earlier parameters
fixed, except for $c_{\beta,m-1}$ when $m>1$, since the preceding
layer's return message first contributes to the loss when layer $m$
is added. We then fine-tune all layers jointly at a reduced learning
rate, with both stages minimizing \eqref{eq:twc-layer-objective}.
The checkpoint minimizing the validation sum $R_H^{(m)}+R_a^{(m)}$,
using the training normalization constants, initializes the next
depth; after depth $L$, the selected parameters are fixed for inference.

\subsection{Network Structure and Inputs}

Each user's real and imaginary antenna maps form $2M$ feature channels
over the $N\times N$ input and output DAFT coordinates, allowing
convolutions to learn dependencies across antennas and coefficients
affected by continuous delays and Doppler shifts. The score network
uses conditional RefineNet \cite{cai2026,song2019}, with the structure
in Table~\ref{tab:twc-networks}. Two fully connected layers embed
$\log(v_\ell)/5$ into 64 features, from which each conditional
normalization layer obtains three coefficients per channel: a weight
for the normalized channel means added to the normalized features,
and a scale and shift for the result.

The classifier combines the denoiser input and conditional channel
estimate into $4M$ channels, scaled by $(1+v_\ell)^{-1/2}$, and
processes them with the convolutions and pooling in
Table~\ref{tab:twc-networks}. Its fully connected layers combine the
pooled features with
\begin{equation}
 \begin{aligned}
 \Big(&\log v_\ell,\
 \log\!\left(1+\frac{\|u_{k\ell}\|^2}{Dv_\ell}\right),\
 \log\!\left(1+\frac{\|\mu_{k\ell}\|^2}{Dv_\ell}\right),\\
 &\log\!\left(1+\frac{\|u_{k\ell}-\mu_{k\ell}\|^2}{Dv_\ell}\right),\
 \logit\rho\Big)
 \end{aligned}
 \label{eq:twc-classifier-features}
\end{equation}
to produce the log-odds correction in \eqref{eq:twc-activity}.
These energy terms describe the input, channel estimate, and residual
relative to the noise level. Since both networks act on individual
users with shared parameters within a layer, their input dimensions
depend on $M$ and $N$, whereas $K$ determines the number of evaluations
and $Q$ enters through the linear estimator.

\begin{table}[t]
\centering
\caption{Network structure for $M=N=16$. Tensor sizes are feature
channels $\times$ height $\times$ width.}
\label{tab:twc-networks}
\footnotesize
\setlength{\tabcolsep}{3pt}
\begin{tabular}{>{\raggedright\arraybackslash}p{0.43\linewidth}
                >{\raggedright\arraybackslash}p{0.49\linewidth}}
\toprule
Stage & Output size or operation\\
\midrule
\multicolumn{2}{l}{\textit{Score network}}\\
Input & $32\times16\times16$; affine map $2u-1$\\
Input convolution, $3\times3$ & $128\times16\times16$\\
Residual groups 1 and 2 & $128\times16\times16$; $256\times8\times8$\\
Residual groups 3 and 4 & $256\times8\times8$ each; dilation 2 and 4\\
Four refinement blocks & Widths 256, 256, 128, 128; final size $16\times16$\\
Output & Conditional normalization, ELU, $3\times3$ convolution; divide by $\sqrt{v/2}$\\
\midrule
\multicolumn{2}{l}{\textit{Activity classifier}}\\
Input & $64\times16\times16$; divide by $\sqrt{1+v}$\\
Three $3\times3$ convolutions & $32\times16\times16$, $64\times8\times8$, $64\times4\times4$\\
Spatial average pooling & 64 features\\
Concatenation & 64 features and the five scalars in \eqref{eq:twc-classifier-features}\\
Two fully connected layers & $69\rightarrow64\rightarrow1$\\
\bottomrule
\end{tabular}
\par\smallskip
\begin{minipage}{0.95\linewidth}
\footnotesize
Each score residual group contains two blocks. Refinement uses
skip connections and bilinear interpolation when sizes differ.
The score network uses exponential linear units (ELUs), while its
noise-conditioning layers and the activity hidden layers use
$x\sigmoid(x)$.
\end{minipage}
\end{table}
